% !TeX program = lualatex
\documentclass[11pt,a4paper]{ltjsarticle}
\usepackage[margin=25mm,headheight=17pt]{geometry}
\usepackage[haranoaji]{luatexja-preset}
\usepackage{amsmath,amssymb,amsthm,mathtools,mathrsfs,booktabs,array,enumitem,longtable}
\usepackage[unicode,hidelinks]{hyperref}
\usepackage{fancyhdr}
\pagestyle{fancy}\fancyhf{}\fancyfoot[C]{\thepage}
\fancyhead[L]{\small\nouppercase{\leftmark}}
\renewcommand{\sectionmark}[1]{\markboth{#1}{}}
\setlength{\emergencystretch}{3em}\allowdisplaybreaks[2]
\setlength{\parskip}{.15em}
\newtheoremstyle{jp}{.6em}{.6em}{\normalfont}{}{\bfseries}{.}{.5em}{}
\theoremstyle{jp}
\newtheorem{theorem}{定理}[section]
\newtheorem{proposition}[theorem]{命題}
\newtheorem{lemma}[theorem]{補題}
\newtheorem{corollary}[theorem]{系}
\newtheorem{definition}[theorem]{定義}
\newtheorem{remark}[theorem]{注意}
\newtheorem{example}[theorem]{例}
\newtheorem{question}[theorem]{問題}
\renewcommand{\proofname}{証明}
\newcommand{\kk}{\mathbb C}\newcommand{\PP}{\mathbb P}
\newcommand{\OO}{\mathcal O}\newcommand{\A}{\mathcal A}
\newcommand{\B}{B_{\mathrm{VK}}}\newcommand{\Blo}{B_{p,q}}
\newcommand{\HH}{\operatorname{HH}}\newcommand{\Hom}{\operatorname{Hom}}
\newcommand{\End}{\operatorname{End}}\newcommand{\Sym}{\operatorname{Sym}}
\newcommand{\Perf}{\operatorname{Perf}}\newcommand{\Der}{\operatorname{Der}}
\newcommand{\Inn}{\operatorname{Inn}}\newcommand{\gr}{\operatorname{gr}}
\newcommand{\rk}{\operatorname{rank}}\newcommand{\Spec}{\operatorname{Spec}}
\newcommand{\id}{\operatorname{id}}\newcommand{\Ext}{\operatorname{Ext}}
\newcommand{\Tor}{\operatorname{Tor}}\newcommand{\Ann}{\operatorname{Ann}}
\newcommand{\coker}{\operatorname{coker}}\newcommand{\Supp}{\operatorname{Supp}}
\newcommand{\pd}{\operatorname{pd}}\newcommand{\depth}{\operatorname{depth}}
\newcommand{\Frac}{\operatorname{Frac}}\newcommand{\Aut}{\operatorname{Aut}}
\newcommand{\Out}{\operatorname{Out}}\newcommand{\GL}{\operatorname{GL}}
\newcommand{\PGL}{\operatorname{PGL}}\newcommand{\Hilb}{\operatorname{Hilb}}
\newcommand{\dashsim}{\mathrel{\overset{\mathrm{bir}}{\sim}}}
\setcounter{tocdepth}{1}

\hypersetup{pdftitle={半直交成分の変形と重み付き射影空間への補完}}
\title{半直交成分の変形と重み付き射影空間への補完\\
\large Veronese--Kronecker代数の二次的変形空間と遅延障害}
\author{}\date{2026年10月29日\quad 日本語公表用原稿VIII・v32}
\begin{document}\maketitle
\begin{abstract}
重み付き射影スタック $\mathbb P(1,1,1,3)$ の標準六項例外列と、次数 $(0,1,3,4)$ の
Veronese--Kronecker部分列を考える。Hochschild一次・二次の制限は十九次元、五十三次元の同型だが、
三次の制限は四十一次元から三十九次元への階数一の写像となる。
六項側の二次障害の線形包は三十八次元で、四項側では全て消える。
四項側のminiversal基底は滑らかな五十三次元形式円板である一方、
六項側は五十三変数と三十八個の同次二次Jacobi方程式で全次数にわたり記述される。
発散零代表によるPoisson構造の持上げ、Kontsevich量子化と二次生成の認定により、
compatibleなhullを閉埋め込みとして構成する。
別稿のPoisson幾何と合わせると、六項側の基底は非被約かつ非純次元である。
さらに二変数の平坦VK族について、標準六項へ補完できる必要十分条件が
任意のArtin基底上で $st=0$ であることを示す。
その帰結として、非零Poisson一次方向を持ちながら、任意に高い次数で初めて補完障害が現れる
全形式変形を得る。別の例外列を許した埋め込みの不存在は主張しない。
\end{abstract}
\tableofcontents
\section{序論と主結果}

半直交部分圏の変形が、もとの周囲の圏の変形から常に誘導されるとは限らない。
本稿ではこの問題を、滑らかな重み付き射影スタック
\[
 X=\mathbb P(1,1,1,3)
\]
の六項例外列と、その四項部分列について具体的に解く。
係数体は全篇を通して $k=\mathbb C$ とする。
重み付き多項式環と二つの有限代数を
\[
 R=k[x_1,x_2,x_3,w],\quad \deg x_i=1,\quad\deg w=3,
 \qquad \Lambda=R_{(0,1,2,3,4,5)},\quad B=R_{(0,1,3,4)}
\]
で定める。次数切り出しの対象 $i<j$ 間の射空間は対応する次数差の $R$ の成分であり、
合成は多項式積である。六項の番号を保つ冪等元
$e=e_0+e_1+e_3+e_4$ に対して $B=e\Lambda e$ となる。
$B$ が本稿でいうVeronese--Kronecker代数である。

Auroux--Katzarkov--Orlovの定理\cite[Theorem 2.12]{AKO}により、
$(\mathcal O_X,\ldots,\mathcal O_X(5))$ はfull strong exceptional collectionであり、
\[
 \Perf(\Lambda)\simeq\Perf(X),\qquad
 \Perf(B)\hookrightarrow\Perf(\Lambda)
\]
は許容部分圏の埋め込みである。
本稿の補完問題は、特殊ファイバー上のこの埋め込みを固定し、
与えられた $B$ の平坦変形を $\Lambda$ の変形のcornerとして実現できるか、という問題である。
頂点冪等元と特殊ファイバーの同定を固定する。
別の例外列や、この特殊ファイバーの埋め込みと結び付かない別の埋め込みについては主張しない。

\subsection{接空間の一致と三次コホモロジーの差}

制限写像を $\rho^i:HH^i(\Lambda)\to HH^i(B)$ と書く。
本稿では次の値を得る。
\[
\begin{array}{c|rrr}
 &HH^1&HH^2&HH^3\\\hline
 \Lambda&19&53&41\\
 B&19&53&39\\
 \operatorname{rank}\rho&19&53&1
\end{array}
\]
特に $\rho^1,\rho^2$ は同型である。
しかし $\ker\rho^3$ は四十次元であり、周囲の三次情報のほとんどが四項側で失われる。
非自明慣性成分を除いた部分には自然な同定
\[
 HH^3(X)_{\rm untw}=H^0(X,\mathcal O_X(6))
 =\operatorname{Sym}^6\langle x_1,x_2,x_3\rangle
 \oplus w\operatorname{Sym}^3\langle x_1,x_2,x_3\rangle\oplus kw^2
\]
があり、この部分の制限の核は最初の二項、すなわち三十八次元である。
これはstacky点 $P=(0:0:0:1)$ で消える反標準切断の空間でもある。

この三十八次元部分は、$HH^2(X)$ の二つの元のSchouten括弧が張る空間に正確に一致する。
したがって、周囲の一次変形に現れる全ての二次障害は四項側で消える。
三十八次元の二次障害の線形包と、四十次元の全制限核は区別する必要がある。

Kuznetsovのnormal Hochschild cohomologyは、半直交分解から生じる制限と変形空間の比較を
一般的に説明する\cite[\S4.2]{Kuznetsov}。
本稿は、二次コホモロジーが一致しても三次で単射性が失われる場合の、具体的な補完軌跡を与える。

\subsection{全次数の変形空間}

$V=H^0(X,\wedge^2T_X)$ と置くと $\dim V=53$ である。
そのJacobi式 $[\pi,\pi]/2=0$ の独立な係数は三十八個の同次二次式
$q_1,\ldots,q_{38}$ となる。
本稿の中心的な結論は、適切な形式座標とcompatibleなminiversal族を選ぶと、corner制限が
\begin{equation}\label{eq:intro-hulls-v31}
 \operatorname{Spf}\frac{k[[u_1,\ldots,u_{53}]]}{(q_1,\ldots,q_{38})}
 \ \hookrightarrow\ \operatorname{Spf}k[[u_1,\ldots,u_{53}]]
\end{equation}
と表示されることである。
左が六項側、右が四項側のminiversal基底である。
これは非標準的なminiversal基底の表示であり、自己同型を持つ変形関手のfine表現ではない。
また三十八は最小二次生成元の数であって、余次元でも正則列の長さでもない。

証明では、重み零多重ベクトルに発散零代表を選び、アフィン空間のKontsevich量子化
\cite{Kontsevich}を全Jacobi軌跡上で同時に適用する。
重み付き発散分解はTang--Wang--Zhang\cite{TWW}、通常の射影空間における対応は
Pym\cite[\S2.2]{Pym}の研究に位置付けられる。
これらの分解や量子化定理自体を新しいものとは扱わない。
ここで必要なのは、それを六項有限代数へ切り出し、二次障害の完全な計算と
Nakayamaの補題を組み合わせて、全次数のminiversal基底を認定することである。

式\eqref{eq:intro-hulls-v31}の重要な帰結として、任意のPoisson一次方向には
補完可能な全形式VK変形が少なくとも一つ存在する。
一方、指定された全形式VK変形の補完可能性は、その一次方向だけでは決まらない。

\subsection{節点による正確な補完条件}

本稿で構成する二変数の階数五十六の平坦族 $B_{s,t}$ については、任意の局所Artin基底 $T$ で
\[
 B_{s,t}\otimes T\text{ が標準六項へ補完可能}
 \quad\Longleftrightarrow\quad st=0\text{ in }T
\]
が成立する。必要性は、特定のOre表示を仮定せず、六項有限代数の全ての変形に対して示す。
証明には $\rho^1$ の全射性と $\rho^2$ の単射性から得られるcorner同型の持上げを使う。

特に $r\ge2$ に対する形式弧 $B_{\hbar,\hbar^r}$ は、非零Poisson一次方向を持つが、
$\hbar^{r+1}$ を法としては補完でき、$\hbar^{r+2}$ を法として初めて補完できなくなる。
従って、全ての形式弧に一様に適用できる固定有限次数のjet判定は存在しない。
この事実は、式 $st=0$ のような有限本の形式方程式による記述と両立する。


\subsection{補完空間の幾何との接続}
Poissonスキーム自体の成分と局所環は、独立した別稿\cite{VKPoisson32}で扱う。
そこでは八つの相異なる成分を認定し、そのうち一つの十九次元成分の一般点に
長さ六の非被約構造を見いだす。本稿のhull同定により、この幾何は六項補完空間にも移る。
これにより、四項側の滑らかさに対して、六項側が非純次元かつ非被約であることが分かる。
非被約構造を捨てたPoisson構造の集合だけでは、Artin基底上の補完問題を記述できない。
本稿では変形理論と補完条件の証明を完結させ、成分ごとの詳細な幾何は上記別稿へ分ける。

\subsection{規約と計算の役割}

多重ベクトルのHKR写像は $p!$ で割った反対称化を用いる。
従って変形乗法 $m+\varepsilon\mu_1$ の半古典括弧は
$\mu_1(f,g)-\mu_1(g,f)$ である。
計算付録の評価行列は整数を保つため階乗を省くが、その階数と像は変わらない。
具体的な二変数族では、この規約により半古典括弧に係数二が現れる。

有限計算は全て有理数の厳密演算で再現できる。
それは接空間、制限写像、独立な障害式を検査するものであり、
全次数の量子化、Artin基底上の必要十分性、形式基底の同型は本文の構造的証明による。

% Standalone section fragment. Requires amsmath, amssymb, amsthm and ordinary
% theorem/proposition/lemma/corollary environments in the parent manuscript.
% All computations concern characteristic zero. Coordinates are numbered 1--3,w.
\section{制限写像の三次までの構造と消える障害}
\label{sec:v31c-restriction}

$V=\langle x_1,x_2,x_3\rangle$、
$R=\operatorname{Sym}(V)[w]$、$\deg x_i=1$、$\deg w=3$ とし、
\[
 X=\mathbb P(1,1,1,3),\qquad
 \Lambda=\operatorname{End}_X\!\left(\bigoplus_{j=0}^{5}\mathcal O_X(j)\right),
 \qquad B=e\Lambda e,\quad e=e_0+e_1+e_3+e_4
\]
と置く。ここで $X$ は滑らかな重み付き射影スタックである。
全対象部分圏への制限により
\[
 \rho^r:\operatorname{HH}^r(X)\longrightarrow\operatorname{HH}^r(B)
\]
を得る。$w$ のみを拡大するトーラスの作用を用い、cochain の重みは
「出力の $w$ 指数から入力の $w$ 指数の和を引いたもの」とする。
以下の同変表示では $V$ を\emph{座標関数}の表現と定める。

\begin{theorem}\label{thm:v31c-restriction}
標数零で次が成り立つ。
\begin{enumerate}
\item $\rho^1$ と $\rho^2$ はそれぞれ $19$ 次元、$53$ 次元の同型である。
\item 全三次制限写像 $\rho^3$ の階数は $1$、核の次元は $40$ である。
その像は $w\partial_{x_1}\wedge\partial_{x_2}\wedge\partial_{x_3}$ の
制限の張る直線である。
\item 恒等慣性成分
$H^0(X,\bigwedge^3T_X)\subset\operatorname{HH}^3(X)$ に制限した核は
$38$ 次元であり、stacky point $P=(0:0:0:1)$ で消える反標準切断に等しい。
\item Schouten bracket の線形像は
\[
 \operatorname{span}\,[\operatorname{HH}^2(X),\operatorname{HH}^2(X)]
 =H^0(X,\mathcal I_P\otimes\mathcal O_X(6))
 \subset H^0(X,\textstyle\bigwedge^3T_X)
\]
である。従ってこの線形像の次元は $38$ であり、
$[\operatorname{HH}^2(B),\operatorname{HH}^2(B)]=0$ である。
\end{enumerate}
\end{theorem}

反標準束と $\mathcal O_X(6)$ の同一視には座標による自明化を用いている。
$\mathrm{GL}(V)$ 同変に記す場合には、後述の行列式指標も付ける。
また、$40$ 次元の全核と $38$ 次元の二次障害の線形像は異なる。
二つの非自明慣性成分を捨てて三次制限写像の全核を記述してはならない。

\subsection{一次制限と有限複体}

重み付きEuler列は
\[
 0\longrightarrow\mathcal O_X\longrightarrow
 \mathcal F=\mathcal O_X(1)^{\oplus3}\oplus\mathcal O_X(3)
 \longrightarrow T_X\longrightarrow0
\]
である。非負次数の線束は高次コホモロジーを持たない。
外積から得られる完全列
$0\to\wedge^{p-1}T_X\to\wedge^p\mathcal F\to\wedge^pT_X\to0$
に帰納法を適用すると、$\wedge^pT_X$ も高次コホモロジーを持たない。
従って恒等慣性成分のHKR群は大域的多重ベクトル場で計算できる。
非自明慣性成分は余次元三の孤立した $B\mu_3$ にあるため、次数二以下には寄与しない
\cite[Corollary 1.4]{FPSS}。


重み零の多項式ベクトル場は
\[
 D=\sum_{i,j=1}^3 a_{ij}x_j\partial_{x_i}
       +(g_3(x)+a w)\partial_w
\]
の形で、$20$ 次元である。Euler 場
\[
 E=\sum_{i=1}^3x_i\partial_{x_i}+3w\partial_w
\]
を割ると $H^0(T_X)$ を得る。
このベクトル場の $B$ への評価が内部導分であると仮定する。
四対象の冪等元に相対的な内部導分は、四頂点のスカラーによるものだけである。
$\operatorname{Hom}(0,1)=V$ 上の作用から $(a_{ij})=\lambda I$ が従う。
$\operatorname{Hom}(0,3)=\operatorname{Sym}^3V\oplus\mathbf k w$ 上では、
純三次式への作用が $3\lambda$ 倍なので、$g_3=0$、$a=3\lambda$ である。
従って $D=\lambda E$ であり、Euler 場の評価は頂点スカラー
$(0,1,3,4)$ による内部導分に一致する。
これで $\rho^1$ の単射性が分かる。

相対正規化 Hochschild 複体は、厳密に増加する対象列と単項式基底から構成する。
$C^0$ は四頂点のスカラーであり、$C^r$ は $r$ 本の非恒等射の可合成列への
多重線形写像の空間である。以下は有理数上の厳密な階数である。
\[
\begin{array}{c|rrrr}
 &C^0&C^1&C^2&C^3\\\hline
\dim&4&620&1584&972
\end{array}
\qquad
\begin{array}{c|rrr}
 &d^0&d^1&d^2\\\hline
\operatorname{rank}&3&598&933
\end{array}
\]
従って $\dim\operatorname{HH}^1(B)=19$ であり、上の単射は同型である。

重み零の多項式 $r$ 重ベクトルを評価する行列を $P_r$ とする。
$r=1,2,3$ の多項式空間の次元はそれぞれ $20,72,92$ であり、
\[
 \operatorname{rank}[d^0\mid P_1]=22,\qquad
 \operatorname{rank}[d^1\mid P_2]=651,\qquad
 \operatorname{rank}[d^2\mid P_3]=934.
\]
第二の等式は $\rho^2$ の像が $651-598=53$ 次元であることを意味する。
重み付き Euler 列から
\[
 h^0(T_X)=20-1=19,\qquad
 h^0(\textstyle\bigwedge^2T_X)=72-19=53
\]
であり、次数二以下には非自明慣性成分の寄与がないので、$\rho^2$ は同型である。
第三の等式は恒等慣性成分からの三次制限の階数が $1$ であることも検証する。
その非零性と核については、次に行列計算によらない証明を与える。

\subsection{恒等慣性成分の三次制限}

多項式三重ベクトルのうち、$\partial_w$ を含むものは $B$ の三次 cochain への
評価が零である。実際、唯一の三本の非恒等射の対象列は
\[
 0\longrightarrow1\longrightarrow3\longrightarrow4
\]
であり、入力の次数は $(1,2,1)$ である。三入力にはいずれも $w$ が現れない。
従って、評価が零とは限らないものは
\[
 F_3(x,w)\partial_{x_1}\wedge\partial_{x_2}\wedge\partial_{x_3},
 \qquad F_3\in R_3
\]
だけである。

Euler 場を含む多重ベクトルが boundary に写ることは、ここでは直接示せる。
$A$ を重み零の多項式 biderivation とし、対象の重みを $h_i$ と書く。
二次 cochain
\[
 f(a:i\to j,b:j\to k)=(h_i+h_j+h_k)A(a,b)
\]
に対して、Leibniz 則を用いると
\[
 (d^2f)(a,b,c)
 =E(a)A(b,c)-E(b)A(a,c)+E(c)A(a,b)
 =(E\wedge A)(a,b,c)
\]
を得る。純三次式 $F_3(x)$ は
$F_3\partial_{x_1}\wedge\partial_{x_2}\wedge\partial_{x_3}=E_x\wedge A$
と書ける。ただし $E_x=\sum x_i\partial_{x_i}$ とする。
$E\wedge A$ との差は $\partial_w$ を含むため、その評価は零である。
従って純三次式を係数とする上の三重ベクトルは boundary に写る。
像は高々
\[
 \tau=w\partial_{x_1}\wedge\partial_{x_2}\wedge\partial_{x_3}
\]
の張る直線である。

\begin{lemma}\label{lem:v31c-functional}
$\tau$ の制限は $\operatorname{HH}^3(B)$ で非零である。
\end{lemma}
\begin{proof}
三次 cochain $\varphi$ に対し、次の整数係数の汎関数を定める。
\[
 \mathcal L(\varphi)=
 \sum_{\sigma\in S_3}\operatorname{sgn}(\sigma)
 \sum_{\ell=1}^{3}
 [x_\ell w]\,
 \varphi(x_{\sigma(1)},x_{\sigma(2)}x_\ell,x_{\sigma(3)}).
\]
各入力には、順に対象 $0\to1\to3\to4$ に対応する始点と終点を付ける。
$[x_\ell w]$ はその出力係数を表す。

任意の二次 cochain $f$ に対して $\mathcal L(d^2f)=0$ である。
微分の四項
$a f(b,c)-f(ab,c)+f(a,bc)-f(a,b)c$ を考える。
第二項は $\sigma(1),\sigma(2)$ の交換で、第三項は
$\sigma(2),\sigma(3)$ の交換で打ち消し合う。
第一項の係数が非零になり得るのは $\ell=\sigma(1)$ のときだけであり、
この場合も $\sigma(1),\sigma(2)$ の交換で打ち消し合う。
第四項は $\ell=\sigma(3)$ に限られ、$\sigma(2),\sigma(3)$ の交換で消える。

非正規化 HKR 代表
\[
 \tau(a,b,c)=w\det\bigl(\partial_{x_i}a,
                         \partial_{x_i}b,
                         \partial_{x_i}c\bigr)
\]
については、一つの $\sigma$ ごとに $\ell$ の和が $1+2+1=4$ となるので
\[
 \mathcal L(\tau)=24\ne0.
\]
標準 HKR 代表の階乗による正規化を用いても非零性は同じである。
\end{proof}

$\Delta=\partial_{x_1}\wedge\partial_{x_2}\wedge
\partial_{x_3}\wedge\partial_w$ と置く。
$E\wedge T=F\Delta$ によって
\[
 H^0(X,\textstyle\bigwedge^3T_X)\simeq R_6
 =\operatorname{Sym}^6V\oplus w\operatorname{Sym}^3V\oplus\mathbf k w^2
\]
と同一視できる。Euler--Koszul 複体の完全性から、この表示は Euler 場を含む
三重ベクトルを正確に割ったものになる。
$E\wedge\tau=-3w^2\Delta$ であり、他の三重ベクトルから $w^2$ 項は生じない。
従って恒等慣性成分での核は
\[
 \operatorname{Sym}^6V\oplus w\operatorname{Sym}^3V
 =\{F\in R_6:[w^2]F=0\}.
\]
これは $P$ で消える切断の空間であり、次元は $28+10=38$ である。

\subsection{非自明慣性成分を含めた全三次制限}

スタック版HKR同型\cite[Corollary 1.4]{FPSS}を用いる。
恒等慣性成分の三次コホモロジーは $39$ 次元である。
これに、孤立した stacky point の二つの非自明慣性成分から、それぞれ
余次元三の normal determinant に対応する一次元の寄与が加わる。
従って
\[
 \dim\operatorname{HH}^3(X)=41.
\]
この三つの正の $w$-重みの類は、$\mathrm{GL}(V)$ 表現として
\[
 \operatorname{HH}^3(X)_{w=1}\simeq\bigl(\det(V)^{-1}\bigr)^{\oplus3}
\]
をなす。恒等慣性成分では $\tau$、非自明成分では局所商
$[\mathbb A^3/\mu_3]$ の normal determinant がこの指標を与える。

一方、$B$ の正の重み部分では $\operatorname{HH}^1$ と
$\operatorname{HH}^2$ は零である。正の cochain は
\[
\begin{aligned}
 C^1_+&\simeq2\operatorname{Sym}^3V^*
                  \oplus V\otimes\operatorname{Sym}^4V^*,\\
 C^2_+&\simeq2V^*\otimes\operatorname{Sym}^2V^*
                  \oplus2V\otimes V^*\otimes\operatorname{Sym}^3V^*,\\
 C^3_+&\simeq V\otimes V^*\otimes\operatorname{Sym}^2V^*\otimes V^*.
\end{aligned}
\]
それぞれの次元は $65,216,162$ である。
この複体の次数一、二での完全性は、二次 cocycle の
\[
 \gamma(v;abz)-\delta(vab;z)=\alpha(v;ab)z-\beta(ab;z)v
\]
を二重に反対称化することから示せる。
反対称部分 $K_{ij,k}=\alpha_{i,jk}-\alpha_{j,ik}$ と
$L_{i,kl}=\beta_{ik,l}-\beta_{il,k}$ は
$2K_{ij,k}+L_{k,ij}=K_{ij,k}+2L_{k,ij}=0$ を満たすため零である。
従って $\alpha,\beta$ は完全対称で、二つの
$\operatorname{Sym}^3V\to\mathbf k$ の boundary により消去できる。
残る $\gamma,\delta$ は四添字で完全対称であり、
$\operatorname{Sym}^4V\to V$ の boundary で消去できる。
一次 cocycle は隣接する射空間で零なので、積による生成から全て零である。

従って Euler character を取ると
\begin{equation}\label{eq:v31c-positive-representation}
 \operatorname{HH}^3(B)_{w=1}
 \simeq \det(V)^{-1}
 \oplus\operatorname{Sym}^3V\otimes\det(V)^{-2}.
\end{equation}
実際、対角トーラスの文字を $t_1,t_2,t_3$ とすれば、その character は
\[
 (t_1t_2t_3)^{-1}
 +(t_1t_2t_3)^{-2}\sum_{a+b+c=3}t_1^at_2^bt_3^c
\]
である。標数零の完全可約性から、これが上の表現分解を与える。

$\rho^3$ は $\mathrm{GL}(V)$ 同変であり、
$\operatorname{Sym}^3V\otimes\det(V)^{-2}$ は
$\det(V)^{-1}$ と同型でない既約表現である。
従って三次の正の重みの像は高々一次元である。
補題\ref{lem:v31c-functional}からその階数は一である。
非正の重みの $38$ 次元は恒等慣性成分に属し、既に零に写ることを示した。
これで全 $\rho^3$ の階数一と核の次元四十が従う。

この議論は非自明慣性成分の個々の HKR 基底の像を決める必要がない。
それらがそれぞれ零に写ると主張しているのではなく、全像の次元を決定している。
HKR 分解の非自明成分への射影により、自然な完全列
\[
 0\longrightarrow H^0(\mathcal I_P\otimes\mathcal O_X(6))
 \longrightarrow\ker\rho^3
 \longrightarrow\operatorname{HH}^3(X)_{\mathrm{twisted}}
 \longrightarrow0
\]
を得る。右の全射性は、恒等慣性成分の $\tau$ が $\rho^3$ の全像を張ることによる。

\subsection{三十八次元の二次障害の線形包}

重み零の多項式二重ベクトルの係数は、$\partial_{x_i}\wedge\partial_{x_j}$
については純二次式、$\partial_{x_i}\wedge\partial_w$ については
純四次式と $w$ 倍の一次式の和である。
従って Schouten bracket を取り、さらに Euler 場と wedge して $R_6$ に
送ったとき、$w^2$ 項は生じない。
よって bracket の線形像は
$\operatorname{Sym}^6V\oplus w\operatorname{Sym}^3V$ に含まれる。

逆の包含には二つの計算で十分である。
\[
 A=x_1^2\partial_{x_1}\wedge\partial_{x_2},\qquad
 B_1=x_2^2\partial_{x_2}\wedge\partial_{x_3},\qquad
 B_2=x_1^4\partial_{x_3}\wedge\partial_w
\]
と置く。座標 Jacobiator の偏極を $J(A,B)$ と書くと、
\[
 E\wedge J(A,B_1)=-6wx_1^2x_2\,\Delta,\qquad
 E\wedge J(A,B_2)=-4x_1^6\,\Delta.
\]
Schouten bracket の標準的な正規化の違いは、ここでは非零の定数倍にすぎない。
$\operatorname{Sym}^3V$ と $\operatorname{Sym}^6V$ は既約であり、bracket は
$\mathrm{GL}(V)$ 同変なので、第一式が $w\operatorname{Sym}^3V$ の全体を、
第二式が $\operatorname{Sym}^6V$ の全体を生成する。
これで三十八次元の像が得られる。

Hochschild cochain の全対象部分圏への制限は Gerstenhaber bracket と可換である。
$\rho^2$ の全射性と、上で得た bracket 像の消滅から
\[
 [\operatorname{HH}^2(B),\operatorname{HH}^2(B)]=0
\]
が従う。これは一次変形の二次障害が全て零であるという主張であり、
それだけから全次数の無障害性を推論しているのではない。
また、三十八個の二次障害式は二次形式のベクトル空間の中で線形独立である。
これを Poisson locus の余次元が三十八であるという主張や、
三十八式が正則列をなすという主張と混同してはならない。

\subsection{再現可能な検証}

\texttt{verify\_ambient\_cohomology\_v31.py} は
$B$ の全相対複体、$20,72,92$ 個の多項式多重ベクトル、
上の十八項汎関数、および全 $\binom{72}{2}$ 個の単項二重ベクトル対を
有理数上で再構成する。三十八個の出力単項式それぞれに対する非零整数係数の
bracket を選び、対角の非零 minor として保存している。

独立の \texttt{verify\_full\_rho3\_v31.py} は六対象の正の重みの複体
\[
 246\longrightarrow1944\longrightarrow4617
 \longrightarrow4374\longrightarrow1458
\]
を再構成する。微分の階数は順に
\[
 246,\quad1698,\quad2916,\quad1458
\]
で、コホモロジーは次数三の三次元に集中する。
$\operatorname{HH}^3(B)_{w=1}$ の双対十一個を六対象の $C^3$ に引き戻し、
$d^3_\Lambda$ の行に付加すると、全階数は一だけ増える。
増加するのは多重次数 $(-1,-1,-1,1)$ の block だけで、そこでの階数は
$294$ から $295$ になる。
これは非自明慣性成分も含む全制限の階数一を、表現論とは別に検証する。
両検証器は標準 Python ライブラリの整数と有理数だけを用い、
通常実行時には certificate を変更せず再計算結果との一致を確認する。

\section{四項側の五十三次元の滑らかな基底}
\label{sec:vk-smooth-hull-v31}

四項側の無障害性は、三次コホモロジーや二次括弧の消滅だけからは従わない。
本節では接写像が全射である滑らかな全形式族を構成する。
これは周囲の二次的基底との比較を自己完結したものにするための補題である。

$K$ を、二頂点の間の射空間が $V_1=k^3$ である三本矢Kronecker代数とする。
四対象を $(0,1)$ と $(3,4)$ の二組に分けると、$B$ は対角ブロックが $K,K$、
貼り合わせ双加群が $M=N\oplus P$ である三角代数となる。
ここで
\[
 N_{ij}=\operatorname{Sym}^{3+j-i}V_1\quad(i,j=0,1),\qquad
 (\dim N_{ij})=\begin{pmatrix}10&15\\6&10\end{pmatrix},
\]
左右の作用は多項式積であり、$P$ は正則 $K$ 双加群で
$\dim(P_{ij})=\left(\begin{smallmatrix}1&3\\0&1\end{smallmatrix}\right)$ である。
$N$ に $w$ 重み零、$P$ に $w$ 重み一を入れる。

\begin{lemma}\label{lem:N-unobstructed-v31}
$N$ の双加群としての変形には、接写像が同型である滑らかな十八次元形式族が存在する。
\end{lemma}
\begin{proof}
まず $A_0=(\operatorname{Sym}V_1)_{(0,1,3,4)}$ に対する相対複体の次元は
$(4,479,1260,810)$、微分の階数は $(3,468,782,0)$ である。
従って $\dim HH^i(A_0)=(1,8,10,28)$ となる。
三次式 $F$ からのJacobian二重ベクトル
\[
 \{f,g\}_F=\det(\nabla f,\nabla g,\nabla F)
\]
の十個の類は $HH^2(A_0)$ の基底をなす。
これらの値と独立性は、本文と同じ相対複体から $w$ を含まない入力・出力を残すことで得られ、
検証器 \texttt{verify\_vk\_smooth\_tangent\_v31.py} が有理数上で確認する。

$\operatorname{End}_{K^e}(N)=k$ である。
実際、四つの成分上の自己写像が左右の多項式積と可換する条件は、全てを同じスカラーにする。
三角代数のHochschild完全列\cite{GG}を使うと、
\[
 \dim\operatorname{Ext}^1_{K^e}(N,N)=18,
 \qquad \dim\operatorname{Ext}^2_{K^e}(N,N)=28.
\]
ここで $HH^1(K)=\mathfrak{sl}(V_1)$、$HH^{\ge2}(K)=0$ であり、
二つの外側の矢の基底変更からの
$\mathfrak{sl}(V_1)^{\oplus2}\to\operatorname{Ext}^1(N,N)$ の核は対角八次元部分である。

Beilinson代数 $D_0=\operatorname{End}_{\mathbb P^2}(\mathcal O\oplus\mathcal O(1)\oplus\mathcal O(2))$
を使う\cite{Beilinson}。
二つの三次元矢空間のテンソル積内の関係三平面を動かすと、
$\operatorname{Gr}(3,9)$ の十八次元の滑らかな形式近傍上に $D_0$ の平坦変形 $D_R$ が得られる。
長さ三の道がないため、関係三平面以外の結合条件はない。

特殊ファイバーのperfect複体 $\mathcal O(3)\oplus\mathcal O(4)$ は自己 $\operatorname{Ext}^1$ と
$\operatorname{Ext}^2$ が零なので、全ての小拡大を通して持ち上がり、同型を除いて一意である。
具体的には有限射影複体の微分を持ち上げ、その二乗の誤差が定める自己 $\operatorname{Ext}^2$ の
classを消して微分を補正する。二つの持上げの差は自己 $\operatorname{Ext}^1$ によって測られる。
両Kroneckerブロック内部と、左の $\mathcal O,\mathcal O(1)$ から
右の $\mathcal O(3),\mathcal O(4)$ への順方向の射複体は、特殊ファイバーで次数零に集中する。
従ってperfect複体に対するNakayamaの補題により、必要な順方向の射空間は
持上げ後も所定の階数の自由加群である。逆方向の射複体に対して同じ集中を仮定しない。
その右側の三次元矢空間のフレームを $\operatorname{PGL}_3$ で独立に動かせば、
滑らかな二十六次元形式基底上に $N$ の変形族を得る。

その接写像は $\operatorname{Ext}^1(N,N)$ に全射である。
上記十個のJacobian括弧による一次積は $D_R$ の一次変形を定め、
次数三・四の加群上でも同じ積を使って対象を持ち上げられる。
持上げの一意性から、それらの接類は $HH^2(A_0)$ の十方向と一致する。
残る八方向は右の矢のフレームの変化が与える。
接写像の核に横断する十八次元形式部分空間に制限すれば、求める族が得られる。
\end{proof}

\begin{theorem}\label{thm:vk-smooth53-v31}
$B$ の標識付き結合代数変形のminiversal基底は
$k[[u_1,\ldots,u_{53}]]$ と同型である。
\end{theorem}
\begin{proof}
$N$ の上記十八次元族と、$P$ の片側を自己同型でねじる八次元族を同時に取る。
正則双加群の長さ一の射影分解をねじると、この二十六次元形式基底上でも
\[
 0\longrightarrow Q_1(P)\xrightarrow{d_P}Q_0(P)\longrightarrow P\longrightarrow0
\]
を得る。$Q_1(P)$ は該当するcornerの射影双加群三個の直和なので、
$\operatorname{Hom}_{K^e}(Q_1(P),N)$ は自由階数 $3\cdot15=45$ である。
任意の $\psi:Q_1(P)\to N$ に対して
\[
 M_\psi=\operatorname{coker}\bigl(Q_1(P)\xrightarrow{(\psi,d_P)}N\oplus Q_0(P)\bigr)
\]
を取ると $0\to N\to M_\psi\to P\to0$ を得る。
$P$ と $N$ が基底上自由なので、$M_\psi$ も自由である。
従って、滑らかな七十一変数の形式基底上に階数五十六の三角代数族が得られる。

特殊点で、その負の重みの接写像は
$\operatorname{Hom}(Q_1(P),N)\twoheadrightarrow\operatorname{Ext}^1(P,N)$ である。
後者は
\[
 [\operatorname{Sym}^3V_1\oplus\operatorname{Sym}^3V_1
 \longrightarrow(\operatorname{Sym}^4V_1)^3],\quad
 (f,g)\longmapsto(x_i(f-g))_i
\]
の余核であり、次元は $45-10=35$ である。
重み零の接像は $N$ と $P$ の変形の全てを含む。
二つの外側のKroneckerブロックの矢の基底変更の像は八次元なので、その次元は
$18+8-8=18$ である。
正の重みの $HH^2(B)$ は零であることが相対複体から分かる。
従って七十一変数族の接写像は、$35+18=53$ 次元の全 $HH^2(B)$ に全射である。

接写像の核に横断する五十三次元形式部分空間に制限する。
miniversal基底を $k[[u_1,\ldots,u_{53}]]/I$ と書くと、この族から形式準同型
$k[[u]]/I\to k[[t_1,\ldots,t_{53}]]$ が得られ、その線形部は可逆である。
形式逆関数定理により合成 $k[[u]]\to k[[t]]$ は同型なので、$I=0$ である。
\end{proof}

本証明は三次コホモロジーが三十九次元存在する場合にも、具体的な全形式族によって
無障害性を示す。後の二変数族の非補完性の証明は、この無障害性を仮定せず成立する。

% Standalone section source for Paper VIII.  Requires amsmath, amssymb, amsthm.
% Statements proved here are independent of the numerical HH restriction test.
\section{重み付き Poisson 構造と外側の量子化}
\label{sec:weighted-poisson-v31}

この節では標数零の体、量子化に関しては $\mathbb C$ 上で考える。
$R=\mathbb C[x_1,x_2,x_3,w]$ に重み $(1,1,1,3)$ を与え、
$X=\mathbb P_{\mathrm{stack}}(1,1,1,3)$ とする。
$E_x=\sum_{i=1}^3x_i\partial_{x_i}$、$E=E_x+3w\partial_w$ と置く。
ここでの量子化の存在は、指定された VK 変形の補完可能性とは異なる。
前者は一次方向を固定して何らかの外側の変形を構成する問題であり、
後者は内側の積の高次項もあらかじめ固定した問題である。

\subsection{Euler 場と発散による一意な持ち上げ}

重み零の多項式 $p$ 重ベクトル場の空間を $\mathcal V^p$ と書く。
標準体積形式
\[
 \Omega=dx_1\wedge dx_2\wedge dx_3\wedge dw
\]
に関する BV 発散を $\delta$ とする。符号は
\[
 \iota_{\delta U}\Omega=(-1)^{p-1}d\iota_U\Omega
 \qquad(U\in\mathcal V^p)
\]
で固定する。特にベクトル場上では通常の発散であり、
二重ベクトル場 $P=\frac12\sum P^{ij}\partial_i\wedge\partial_j$
上では $\delta P=\sum_{i,j}\partial_i(P^{ij})\partial_j$ である。
$\delta^2=0$ であり、$\delta$ は Schouten 括弧を生成する BV 作用素である。

\begin{lemma}[重み付き Euler ホモトピー]
\label{lem:weighted-euler-homotopy-v31}
任意の重み零の多項式多重ベクトル場 $U$ に対して
\[
 \delta(E\wedge U)+E\wedge\delta U=6U
\]
が成り立つ。したがって
\[
 L(U):=U-\frac16E\wedge\delta U
\]
は $\ker\delta$ への射影であり、Euler 場による商から同型
\[
 \mathcal V^p/(E\wedge\mathcal V^{p-1})
 \xrightarrow{\ \sim\ }
 \ker(\delta:\mathcal V^p\longrightarrow\mathcal V^{p-1})
\]
を誘導する。
\end{lemma}

\begin{proof}
一般の多重ベクトル場について BV 恒等式は
\[
 \delta(E\wedge U)+E\wedge\delta U
 =(\operatorname{div}E)U+[E,U]
\]
を与える。ここでは $\operatorname{div}E=6$、$[E,U]=0$ である。
これと $\delta^2=0$ から $\delta L(U)=0$、$L^2=L$ が従う。
$L(E\wedge V)=0$ であり、$U-L(U)$ は Euler 場を因子に持つ。
これで存在と一意性がともに示された。
より直接には、$E\wedge U=0$ かつ $\delta U=0$ なら、
上の恒等式を $U$ に適用して $6U=0$ を得る。
\end{proof}

重み付き Euler 完全列と正の線束の高次コホモロジーの消滅から
\[
 H^0(X,\wedge^pT_X)
 =\mathcal V^p/(E\wedge\mathcal V^{p-1})
\]
となる。この同定は Schouten 括弧を保つ。
実際、重み零の場に対する Euler 場を因子に持つ場は Lie イデアルであり、
$\ker\delta$ は BV 恒等式により Schouten 括弧で閉じている。

\begin{proposition}[Poisson 構造の一意な unimodular lift]
\label{prop:weighted-poisson-lift-v31}
$X$ の大域的二重ベクトル場 $\pi$ を、多項式二重ベクトル場 $P$ で表す。
このとき
\[
 \widehat\pi=P-\frac16E\wedge\delta P
\]
は代表元によらず、$\delta\widehat\pi=0$ を満たす唯一の重み零代表である。
さらに $\pi$ が Poisson であることと、
$\widehat\pi$ が $\mathbb A^4$ 上の通常の多項式 Poisson 構造であることは同値である。
\end{proposition}

\begin{proof}
唯一性は補題から従う。$\pi$ が Poisson なら
$R_3=[\widehat\pi,\widehat\pi]$ は射影化後に零になるので
$E\wedge R_3=0$ である。一方、BV 恒等式と
$\delta\widehat\pi=0$ から $\delta R_3=0$ である。
補題の最後の観察により $R_3=0$ となる。逆向きは射影化すればよい。
\end{proof}

同じ証明は正の重み $(a_0,\ldots,a_n)$ に対し、$6$ を
$\sum a_i$ に置き換えて成り立つ。
通常の射影空間における unimodular lift は
\cite[\S 2.2]{Pym} に現れ、重み付き多項式 Poisson 構造の対応する分解は
\cite{TWW} で扱われている。
上の証明は本稿に必要な射影商からの持ち上げを明示したものである。

\begin{corollary}[すべての Poisson 一次方向の積分]
\label{cor:all-poisson-directions-v31}
$\pi\in H^0(X,\wedge^2T_X)$ が Poisson なら、
$\pi$ を一次方向とする重み付き次数を保つ平坦形式量子化 $S_\hbar$ が存在する。
したがって六対象の代数
\[
 \Lambda_\hbar
 =\bigoplus_{0\leq i\leq j\leq5}(S_\hbar)_{j-i}
\]
は $\Lambda$ の形式変形となる。各次数成分の階数は可換な場合と等しく、
$\Lambda_\hbar$ の階数は $141$ である。
この主張は任意の完備局所基底上の形式 Poisson テンソルにも適用できる。
\end{corollary}

\begin{proof}
命題により $\pi$ を重み零の多項式 Poisson 構造
$\widehat\pi$ に持ち上げ、アフィン空間上の Kontsevich 量子化
\cite{Kontsevich} を適用する。
各グラフ演算は線形座標変換に同変である。
重み付き $\mathbb G_m$ 作用は線形作用なので、各次数の積は重みを保つ。
したがって $(S_\hbar)_d=R_d[[\hbar]]$ を台とする積が得られる。
第一項の規約は採用した HKR 同定と一致させる。
この積を次数差 $0,\ldots,5$ に評価すれば、表示した有限自由代数を得る。
結合則は $S_\hbar$ の結合則そのものである。
\[
 (\dim R,\ldots,\dim R_5)=(1,3,6,11,18,27)
\]
なので階数は
$6+5\cdot3+4\cdot6+3\cdot11+2\cdot18+27=141$ である。
一般の完備局所基底では Poisson テンソルの係数を極大イデアルに取れば、
同じグラフ級数がそのイデアルの位相で収束する。
\end{proof}

この証明だけから一般の量子化の AS 正則性、代数化、収束性を結論してはいけない。
ここで必要な結論は、標準六対象の有限代数が形式的に平坦に変形されることである。
また、この系は同じ一次方向を持つすべての $B_\hbar$ の補完を主張しない。

\subsection{53 変数における明示的な Poisson 方程式}

$S_d=\mathbb C[x_1,x_2,x_3]_d$ と書く。
一般の重み零二重ベクトル場は
\[
 P=\Pi(x)+(F(x)+wA(x))\wedge\partial_w,
 \qquad F\in S_4^{\oplus3},\quad A\in S_1^{\oplus3}
\]
と書ける。Euler 場の外積を足して $wA$ を消すと、
\[
 H^0(X,\wedge^2T_X)
 \simeq (S_2^{\oplus3})\oplus(S_4^{\oplus3}/S_3x)
\]
となる。実際、
\[
 E\wedge\bigl(a+(g+bw)\partial_w\bigr)
 =E_x\wedge a+\{gx+w(bx-3a)\}\wedge\partial_w
\]
において $a=(A+bx)/3$ とすればよい。
$w$ に依存しない代表の間に残る同値は $F\mapsto F+gx$ のみである。
次元は $18+(45-10)=53$ である。

\[
 \Pi=p_1\partial_2\wedge\partial_3
       +p_2\partial_3\wedge\partial_1
       +p_3\partial_1\wedge\partial_2
\]
と書き、
\[
 c=x\cdot p\in S_3,\qquad D=\operatorname{curl}p\in\mathfrak{sl}_3
\]
と置く。ここで $\mathfrak{sl}_3$ は発散零の線形ベクトル場として作用する。
ベクトル解析の恒等式から
\[
 p=\frac13(\nabla c-x\times D)
\]
であり、これは $S_2^{\oplus3}\simeq S_3\oplus\mathfrak{sl}_3$ を与える。
三変数の二次 Poisson 構造に関するこの分解は古典的である
\cite{LiuXu1992,DoninMakarLimanov1998}。

\begin{proposition}[明示的な二次方程式]
\label{prop:weighted-poisson-equations-v31}
上の標準形 $P=\Pi+F\wedge\partial_w$ が $X$ 上で Poisson となる条件は
\begin{align}
 D(c)&=0,\label{eq:weighted-base-jacobi-v31}\\
 2F(c)-(x\times D)\cdot F-c\operatorname{div}F&=0.
 \label{eq:weighted-mixed-jacobi-v31}
\end{align}
第一式は三次式の $10$ 係数、第二式は六次式の $28$ 係数を与える。
これらは $53$ 個の座標に関する $38$ 個の一次独立な二次式である。
\end{proposition}

\begin{proof}
$[P,P]$ の $\partial_1\wedge\partial_2\wedge\partial_3$ 成分は
$[\Pi,\Pi]$ であり、残りは $2[\Pi,F]\wedge\partial_w$ である。
$E\wedge[P,P]=0$ において $w$ の係数と定数項を分けると
$[\Pi,\Pi]=0$ と $E_x\wedge[\Pi,F]=0$ を得る。
前者は $p\cdot\operatorname{curl}p=0$、したがって $D(c)=0$ と同値である。

$[F,\Pi]$ の係数ベクトルを $q$ とすると
\[
 q=F(p)+(\operatorname{Jac}F)^tp-(\operatorname{div}F)p.
\]
$F$ の各係数は四次なので
\[
 x\cdot q=F(c)+3p\cdot F-c\operatorname{div}F
 =2F(c)-(x\times D)\cdot F-c\operatorname{div}F.
\]
$F\mapsto F+gx$ による不変性は
$E_x(c)=3c$、$\operatorname{div}(gx)=6g$ から従う。

最後に二次式の一次独立性を示す。
$F=0$ と置くと第一式のみが残り、
$\mathfrak{sl}_3S_3=S_3$ によりその $10$ 係数は独立である。
次に $c=0$ と置く。$D$ を非対角行列単位に取ると、
$(x\times D)_i$ のいずれかを $\pm x_j^2$ とできる。
任意の六次単項式は少なくとも一つの $x_j^2$ で割り切れるので、
$F_i$ を適切な四次単項式とすれば第二式は任意の六次単項式を張る。
よって第二式の $28$ 係数も独立である。
\end{proof}

この標準形に対する unimodular lift は明示的に
\[
 \widehat P
 =\Pi+\frac16E_x\wedge D
 +\left(F-\frac{\operatorname{div}F}{6}x-\frac w2D\right)
       \wedge\partial_w
\]
である。方程式が $38$ 個あることは正則列であることを意味しない。
成分とその非被約構造の幾何は別稿\cite{VKPoisson32}で独立に証明する。


\section{周囲の変形の二次的なminiversal基底}
\label{sec:quadratic-hull-v31}

本節の量子化には複素数体を用いる。
以下の結論はminiversal基底の存在とその形式同型型について述べるものであり、
自己同型を無視したfine moduli空間の存在や、座標の標準性を主張しない。

\subsection{二次的な基底を認定する補題}

\begin{lemma}[二次生成と量子化族による認定]
\label{lem:quadratic-hull-recognition-v31}
形式変形問題が有限次元の接空間 $V$ とminiversal基底
$R=P/I$、$P=k[[u_1,\ldots,u_d]]$、$I\subset\mathfrak m_P^2$ を持つとする。
同次二次式で生成されるイデアル $J\subset P$ について、次を仮定する。
\begin{enumerate}
\item $P/J$ 上に、接写像が $V$ の同型となる全形式変形が存在する。
\item この接空間の同定の下で、miniversal変形の二次障害方程式は
      $J_2\subset\operatorname{Sym}^2(V^*)$ と一致する。
\end{enumerate}
このとき $R$ は $P/J$ と形式的に同型である。
\end{lemma}

\begin{proof}
形式的versalityを各Artin商に逐次適用して、与えられた全形式変形を誘導する
連続な局所準同型 $\phi:R\longrightarrow P/J$ を得る。
その接写像は同型なので、$P\to R\xrightarrow{\phi}P/J$ の各変数の像を
$P$ に持ち上げると、可逆な一次項を持つ $P$ の形式自己同型を得る。
この自己同型で座標を変えれば
\[
 I\subset J,\qquad I_2=J_2
\]
とできる。ここで $I_2=(I+\mathfrak m_P^3)/\mathfrak m_P^3$ である。

$J$ は同次二次式で生成されるから
\[
 J\cap\mathfrak m_P^3=\mathfrak m_PJ.
\]
各二次生成元 $q$ に対し、$q$ と同じ二次項を持つ $f\in I$ を選ぶと、
$q-f\in J\cap\mathfrak m_P^3=\mathfrak m_PJ$ である。
従って $J=I+\mathfrak m_PJ$ となる。
有限生成 $P$ 加群 $J/I$ にNakayamaの補題を適用すれば $J/I=0$、すなわち $I=J$ である。
\end{proof}

この補題は全ての高次障害を個別に計算する必要を取り除く。
二次の必要条件を満たす全ての点を一つの全形式族で実現できれば、
それ以外の独立な高次方程式は残らない。

前節の発散零代表により、大域的多重ベクトルのSchouten代数を
重み零多項式多重ベクトルのSchouten代数へ持ち上げる切断を得ている。

\subsection{全Jacobi軌跡の量子化}

\[
 V=H^0(X,\wedge^2T_X),\qquad
 W=\operatorname{im}\bigl(\operatorname{Sym}^2V
       \xrightarrow{[\ ,\ ]}H^0(X,\wedge^3T_X)\bigr)
\]
と置く。既に示した計算により $\dim V=53$、$\dim W=38$ である。
$V$ の基底の発散零代表を $\Pi_1,\ldots,\Pi_{53}$ とし、
\[
 P=\mathbf C[[u_1,\ldots,u_{53}]],\qquad
 \Pi=\sum_{a=1}^{53}u_a\Pi_a
\]
と置く。
Jacobi式 $\frac12[\Pi,\Pi]$ の $W$ の基底に関する係数を
$q_1,\ldots,q_{38}$ とし、$J=(q_1,\ldots,q_{38})$ とする。
これらは線形独立な同次二次式である。

$P/J$ 上では $\Pi$ がアフィン空間 $\mathbf A^4$ 上のPoisson構造を定める。
Kontsevichのアフィン量子化公式\cite[\S6.3--6.4]{Kontsevich}をパラメータの極大イデアルに関して適用し、
\[
 f\star g=fg+\sum_{n\ge1}\frac1{n!}U_n(\Pi,\ldots,\Pi)(f,g)
\]
を得る。
各有限パラメータ次数では有限個のgraph演算だけが現れるので、これは
$P/J$ の極大イデアルに関するadic位相で良く定義された結合的な乗法である。
graph演算は多項式係数を保ち、線形座標変換に自然である。
特に重み付き $\mathbf G_m$ 作用を保つため、$R$ の重み付き次数も保つ。
従って各次数成分は元の $R$ の次数成分の自由な形式変形である。
これを六項の次数 $(0,1,2,3,4,5)$ に切り出せば、
$P/J$ 上の階数百四十一の全形式変形 $\Lambda_{\mathrm{K}}$ が得られる。

Kontsevichの $U_1$ は通常のHKR反対称化
\[
 X_1\wedge\cdots\wedge X_m
 \longmapsto\frac1{m!}\sum_{\sigma\in S_m}
       \operatorname{sgn}(\sigma)X_{\sigma(1)}\otimes\cdots\otimes X_{\sigma(m)}
\]
である。このため、六項の有限代数へ評価した一次写像の四項cornerへの制限は、
既に検証した二重ベクトルの評価行列と同じ写像である。
反対称化を階乗無しで定義した計算規約とは、二重ベクトル上で非零定数倍だけ異なる。
従ってその四項への制限 $V\to HH^2(B)$ は同型であり、
$HH^2(\Lambda)\to HH^2(B)$ の同型から、
$\Lambda_{\mathrm K}$ の接写像 $V\to HH^2(\Lambda)$ も同型である。

\begin{theorem}[周囲のminiversal基底の二次表示]
\label{thm:ambient-quadratic-hull-v31}
$\Lambda$ の標識付き結合代数変形のminiversal基底は、形式座標変換を除き、
\[
 \mathbf C[[u_1,\ldots,u_{53}]]/(q_1,\ldots,q_{38})
\]
である。ここで $q_a$ は大域的な重み零二重ベクトルのJacobi式の独立な係数である。
\end{theorem}

\begin{proof}
上の量子化族は、全Jacobi軌跡上の全形式変形であり、接写像は同型である。
また、Hochschild二次障害はGerstenhaber括弧 $\frac12[\xi,\xi]$ である。
既に得たHochschild--HKRの同定とSchouten括弧の計算から、
その独立な二次成分は正確に $q_1,\ldots,q_{38}$ である。
従って補題\ref{lem:quadratic-hull-recognition-v31}が適用できる。
\end{proof}

本定理で三十八は最小二次生成元の個数を意味する。
方程式間のsyzygy、成分の次元、既約性、完全交叉性は、この二次表示だけからは決まらない。
非純次元性・非被約性・非完全交叉性は、節~\ref{sec:ambient-geometry-consequences-v32}で別稿の幾何から導く。
特に、この表示だけから次元が $53-38=15$ であるとは結論できない。

\subsection{適合するminiversal基底上の閉埋め込み}

\begin{proposition}[corner写像の基底上の表示]
\label{prop:compatible-hull-closed-embedding-v31}
$B$ の変形が無障害であり、そのminiversal基底が
$P_B=\mathbf C[[v_1,\ldots,v_{53}]]$ であるとする。
corner制限 $HH^2(\Lambda)\to HH^2(B)$ が同型であるとき、
corner制限に適合するminiversal基底間の写像を選ぶことができ、その写像は閉埋め込み
\[
 \operatorname{Spf}R_\Lambda\hookrightarrow\operatorname{Spf}P_B
\]
である。適切な形式座標の下で、これは前節のJacobi方程式 $J$ による閉埋め込みとなる。
\end{proposition}

\begin{proof}
$R_\Lambda=P_u/J$ 上のminiversalな周囲の変形のcornerを取る。
$B$ の形式的versalityを各Artin商に逐次適用すれば、そのcornerを誘導する
連続準同型 $f:P_B\to R_\Lambda$ を得る。
接写像は $HH^2$ の制限写像であるから同型である。
各 $f(v_i)$ を $P_u$ の元 $F_i(u)$ に持ち上げると、一次行列が可逆な写像
\[
 F:P_B\xrightarrow{\ \sim\ }P_u
\]
を得る。形式逆写像定理により $F$ は同型であり、$f$ は $F$ と商写像
$P_u\twoheadrightarrow P_u/J$ の合成である。
従って $f$ は全射で、核は $F^{-1}(J)$ である。
\end{proof}

この表示は選択したminiversal族と形式座標に依存する。
変形groupoidのfine代表性やgroupoidそのものの閉埋め込みを主張するものではない。
また $\rho^3$ 全体の核と、一次変形の二次障害の線形包とは区別する。
本例では前者は四十次元、後者は三十八次元であり、
二次方程式の個数を定めるのは後者である。


\section{二変数族の補完可能性と任意に遅い障害}
\label{sec:completion-locus-v31}

本節では標数零の体 $k$ 上で考える。
重み付き多項式環 $R=k[x_1,x_2,x_3,w]$ に
$\deg x_i=1$、$\deg w=3$ を入れ、
\[
 \Lambda=(R)_{(0,1,2,3,4,5)},\qquad
 e=e_0+e_1+e_3+e_4,\qquad B=e\Lambda e
\]
と置く。有限代数の次数切り出しは、対象 $i<j$ 間の射空間を
$(R)_{j-i}$ とし、合成を環の乗法から定めるものを意味する。
射の積の順序は本節を通して同じ規約を用いる。
以下の変形と同型は、特殊ファイバーの標識と各頂点冪等元を保つ。

\begin{definition}[標準六項への補完]
局所Artin $k$ 代数 $T$ 上の $B$ の変形 $B_T$ が
\emph{標準六項へ補完可能}であるとは、
$\Lambda$ の平坦変形 $\Lambda_T$ と標識に整合する同型
$e\Lambda_Te\simeq B_T$ が存在することをいう。
\end{definition}

これは固定した四項部分列と六項列の組を変形する問題である。
中央ファイバーの埋め込みと結び付かない別の例外列や、
別の非可換重み付き射影空間への埋め込みを排除する定義ではない。

\subsection{接空間の単射性から得る補完の剛性}

まず、具体例から独立な補題を用意する。

\begin{lemma}[corner同型の持上げ]
\label{lem:corner-isomorphism-lifting-v31}
$\Lambda$ を有限次元 $k$ 代数、$B=e\Lambda e$ とする。
頂点に関する相対Hochschild複体の制限が
\[
 HH^1(\Lambda)\longrightarrow HH^1(B)\quad\text{で全射},
 \qquad
 HH^2(\Lambda)\longrightarrow HH^2(B)\quad\text{で単射}
\]
であると仮定する。
任意の局所Artin基底 $T$ 上で、$\Lambda$ の二つの平坦変形
$\Lambda_T,\Lambda'_T$ のcorner間の標識付き同型は、
周囲の代数間の標識付き同型へ持ち上がる。
特に、同じcornerを持つ周囲の変形は同型を除き高々一つである。
同型そのものの一意性は主張しない。
\end{lemma}

\begin{proof}
$T$ の長さについて帰納する。
$T\twoheadrightarrow \overline T$ を、核 $I$ が
$\mathfrak m_T I=0$ を満たす小拡大とする。
帰納法により、与えられたcorner同型の $\overline T$ への還元は
周囲の代数同型に持ち上がる。
これを頂点成分ごとに $T$ 上の線形同型へ持ち上げる。
corner上では最初に指定した同型そのものを使える。
有限平坦加群は自由なので、この線形同型は可逆である。

この線形同型で乗法を移送すると、二つの乗法の差は
\[
 c\in Z^2(\Lambda)\otimes_k I,
 \qquad \operatorname{res}(c)=0
\]
となる。$HH^2$ 上の単射性から $c=d\xi$ と書ける。
このとき $\operatorname{res}(\xi)$ は $B$ の一コサイクルである。
$HH^1$ 上の全射性により、周囲の一コサイクル $\eta$ を選び、
$\operatorname{res}(\xi)-\operatorname{res}(\eta)=db$ とできる。
cornerの零コチェイン $b$ は周囲の零コチェインへ持ち上がるため、
内的導分を加えて $\operatorname{res}(\eta)=\operatorname{res}(\xi)$ とできる。
従って $\xi-\eta$ はcorner上零で、$d(\xi-\eta)=c$ を満たす。
この無限小の基底変更で乗法の差を消せば、指定したcorner同型を保ったまま
周囲の代数同型が得られる。
\end{proof}

定理\ref{thm:v31c-restriction}の $HH^1,HH^2$ の同型により、
補題\ref{lem:corner-isomorphism-lifting-v31}をこの組に適用できる。
この帰結に、四項側の全変形関手の無障害性は必要ない。

\subsection{四項側の全二変数変形}

$\mathscr R=k[[s,t]]$ と置き、
\[
 A_s=
 \frac{\mathscr R\langle x_1,x_2,x_3\rangle}
 {([x_1,x_2]-2s x_1x_3,\,[x_1,x_3],\,[x_2,x_3])}
\]
を取る。
これは $\mathscr R[x_1,x_3]$ のOre拡大であり、順序付き単項式が基底となる。
その四項切り出しを $N_s=(A_s)_{(0,1,3,4)}$ とする。
四頂点を順に $0,1,2,3$ と番号付けた記号では、自由加群 $M_t$ を
\[
 (M_t)_{02}=\mathscr Rw_0,\qquad
 (M_t)_{13}=\mathscr Rw_1,\qquad
 (M_t)_{03}=\bigoplus_{i=1}^3\mathscr Rv_i
\]
で定める。非恒等射による非零の作用は
\begin{equation}\label{eq:two-parameter-module-v31}
 x_iw_1=(1+t\delta_{i2})v_i,
 \qquad w_0x_i=(1-t\delta_{i2})v_i
\end{equation}
だけとする。
\[
 B_{s,t}=N_s\ltimes M_t
\]
と定める。$M_t$ の射は少なくとも一つの中間頂点を飛ばすので、
$M_t$ の元を含む三本の非恒等射の可合成列は存在しない。
従って\eqref{eq:two-parameter-module-v31}に追加の三重合成条件はなく、
$B_{s,t}$ は結合的である。階数は $51+5=56$ であり、
$\mathscr R$ 上自由な $B$ の全形式変形を与える。

\begin{theorem}[補完可能性を表す節点]
\label{thm:exact-node-completion-v31}
任意の局所Artin $k$ 代数 $T$ と $s,t\in\mathfrak m_T$ に対し、
次は同値である。
\begin{enumerate}
\item $B_{s,t}\otimes_{\mathscr R}T$ が標準六項へ補完可能である。
\item $st=0$ が $T$ 内で成立する。
\end{enumerate}
従って、この二変数族における補完の存在条件は形式的節点
\[
 \operatorname{Spf}k[[s,t]]/(st)
 \ \subset\ \operatorname{Spf}k[[s,t]]
\]
によって、全てのArtin基底上で表される。
これは補完の同型全体を表すfine moduli空間という主張ではなく、
指定した二変数族内の補完存在条件の表示である。
\end{theorem}

\begin{proof}[十分性]
$st=0$ とする。$1\pm t$ は単元なので
\[
 \lambda(t)=\frac{1-t}{1+t}
\]
を定められる。$\sigma_t(x_1)=x_1$、$\sigma_t(x_3)=x_3$、
$\sigma_t(x_2)=\lambda(t)x_2$ とすると
\[
 \sigma_t\bigl([x_1,x_2]-2s x_1x_3\bigr)
 =2s\bigl(\lambda(t)-1\bigr)x_1x_3
 =-\frac{4st}{1+t}x_1x_3=0.
\]
従って $\sigma_t$ は $A_s$ の次数を保つ自己同型である。
重み三の変数を加えたOre拡大
\[
 S_{s,t}=A_s[w;\sigma_t],\qquad wa=\sigma_t(a)w
\]
は平坦であり、順序付き単項式 $x_1^ax_2^bx_3^cw^d$ を基底に持つ。
その六項切り出し $\Lambda_{s,t}=(S_{s,t})_{(0,1,2,3,4,5)}$ は
階数百四十一の平坦な周囲の変形である。
cornerでは $w_0,w_1$ をそれぞれの次数三の $w$ とし、
\[
 v_i=\frac{x_iw}{1+t\delta_{i2}}
\]
と取れば、積は正確に\eqref{eq:two-parameter-module-v31}となる。
従って $e\Lambda_{s,t}e\simeq B_{s,t}$ である。
\end{proof}

\begin{proof}[必要性]
二つの軸上の周囲の変形の一次classを
$\alpha_s,\alpha_t\in HH^2(\Lambda)$ と書く。
四項側の一次積を重み零の反対称二コチェインと比べると、
$s$ 軸では $D=x_1x_3\partial_{x_1}\partial_{x_2}$ のHochschild境界だけ異なり、
$t$ 軸ではそのまま一致する。
従って $HH^2$ の制限の単射性により、周囲の二つのclassも対応する大域二重ベクトルのclassである。
以下では一次交換子を半古典括弧とする規約を使う。
これらの混合障害は非零である。
実際、$x_3\ne0$ のチャートで
\[
 x=x_1/x_3,\quad y=x_2/x_3,\quad z=w/x_3^3
\]
とすると、二つの半古典括弧は
\[
 \{x,y\}_s=2x,\qquad \{y,z\}_t=2yz
\]
を与える。ほかの生成元対の括弧は零である。
それらの和のJacobi式の混合項は $4xz$ で非零だから、
HKRの括弧との整合性により
\begin{equation}\label{eq:mixed-obstruction-nonzero-v31}
 [\alpha_s,\alpha_t]\ne0\quad\text{in }HH^3(\Lambda).
\end{equation}
この非消滅は滑らかな通常のアフィンチャートで検出されるため、
stacky点のHochschild寄与の扱いに依存しない。

いま $T$ 内で $st\ne0$ であると仮定する。
$J=T(st)$ とし、
\[
 T'=T/\mathfrak m_TJ,\qquad I=J/\mathfrak m_TJ
\]
と置く。Nakayamaの補題により $I\ne0$ であり、
$\mathfrak m_{T'}I=0$ を満たす。
従って $T'\twoheadrightarrow\overline T=T'/I$ は小拡大であり、
$\overline T$ 上では $st=0$ が成立する。

もし $T$ 上に補完があれば、その $T'$ への基底変換も補完である。
その $\overline T$ への還元は、十分性で構成した標準Ore補完と同じcornerを持つ。
補題\ref{lem:corner-isomorphism-lifting-v31}により、この二つの周囲の変形は同型である。
従って、標準Ore補完の $\overline T$ から $T'$ への障害を計算すればよい。

標準Ore補完を $k[[s,t]]/(st)$ 上の固定PBW基底で記述する。
その乗法は一意に
\[
 \mu=\mu_0+\sum_{i\ge1}s^i\alpha_i+\sum_{j\ge1}t^j\beta_j
\]
と書ける。各軸上では全次数で結合的であり、
$[\alpha_1]=\alpha_s$、$[\beta_1]=\alpha_t$ である。
同じ式を $T'$ 上へ持ち上げた線形乗法を $\widetilde\mu$ とする。
その結合子の純粋な $s$ 項・$t$ 項は零である。
混合項 $s^it^j$ のうち $i+j\ge3$ のものも
$\mathfrak m_{T'}I=0$ によって消える。
従って結合子が与える障害classは
\[
 (st)\,[\alpha_s,\alpha_t]\in HH^3(\Lambda)\otimes_k I
\]
である。$st$ は $I$ の非零生成元であり、
\eqref{eq:mixed-obstruction-nonzero-v31}からこのclassは非零である。
別の線形乗法の持上げに変えても変化はHochschild coboundaryだけなので、
結合的な持上げは存在しない。これは補完の存在に矛盾する。
\end{proof}

\subsection{非零Poisson一次方向と任意に遅い非延長}

\begin{corollary}[任意に高い次数で最初に現れる補完障害]
\label{cor:arbitrarily-delayed-completion-v31}
整数 $r\ge2$ に対し
\[
 B^{(r)}_\hbar=B_{\hbar,\hbar^r}
\]
と置く。これは $k[[\hbar]]$ 上の階数五十六の平坦変形である。
その一次方向は非零Poisson方向であり、同じ一次方向は周囲でも全次数に積分できる。
しかし $B^{(r)}_\hbar$ 自身は、標準六項へ形式的に補完できない。
より正確には
\[
 B^{(r)}_\hbar\bmod\hbar^{r+1}
 \quad\text{は補完可能},\qquad
 B^{(r)}_\hbar\bmod\hbar^{r+2}
 \quad\text{は補完不能}
\]
である。
\end{corollary}

\begin{proof}
一次方向は $s$ 軸の方向であり、$t=0$ のOre拡大が全次数の周囲の積分を与える。
一方、$st=\hbar^{r+1}$ なので、定理\ref{thm:exact-node-completion-v31}を
二つのArtin商へ適用すれば結論を得る。
\end{proof}

従って、与えられた全形式VK変形の補完可能性を、その一次方向のPoisson性だけで判定することはできない。
さらに、任意の整数 $n\ge1$ に対し、補完可能な $B_{\hbar,0}$ と
補完不能な $B_{\hbar,\hbar^{n+1}}$ は $\hbar^{n+1}$ を法として一致する。
従って、全ての形式弧に一様に適用できる有限 $\hbar$-adic jetによる判定法は存在しない。
これは形式的な補完軌跡に有限本の方程式が存在しないという意味ではない。
実際、本節の二変数族では、その全方程式は一つの式 $st=0$ である。

この族の接平面内では、半古典二重ベクトルのPoisson条件も正確に $st=0$ である。
この意味で、この二変数族では二次のPoisson方程式が全補完条件を尽くす。
任意の五十三次元VK変形について、一次Poisson方向が与えられた全形式変形の補完を保証するという主張とは区別する。

\section{補完空間の非被約性と非純次元性}
\label{sec:ambient-geometry-consequences-v32}

前節までのhull同定は、Poissonスキームの幾何を変形空間へ移す。
ここでは別稿\cite{VKPoisson32}の結果を用い、その変形論的帰結を明記する。
成分と形式局所環の証明は、そこで本稿のhull同定とは独立に与えられる。

\begin{corollary}[非被約な六項補完空間]
\label{cor:ambient-nonreduced-v32}
六項代数 $\Lambda$ のminiversal hullは非被約かつ非純次元であり、完全交叉ではない。
対応する代数的Jacobiスキームは、次元
\[
 35,\ 28,\ 23,\ 22,\ 19,\ 19,\ 19,\ 19
\]
の相異なる既約成分を少なくとも持つ。
最後の十九次元成分の一般点では局所環の長さは六である。
これらの次元はスカラー倍や座標自己同型で商する前の次元である。
\end{corollary}

\begin{proof}
別稿の成分定理と共点三直線の局所模型による。
特に、ある点の完備局所環は滑らかな十八変数を除いて
\[
 \mathbb C[[a,b,u]]/(ab,b^2,a^5u)
\]
となり、$u$ を可逆にすると長さ六の非被約環を得る。
有限型の $\mathbb C$ 代数はexcellentであり、被約局所環の完備化は被約なので、
この局所模型は元のJacobiスキームの非被約性を意味する。
Jacobiイデアルは同次である。非被約軌跡は閉かつスカラー作用で不変であり、
その閉包は原点を含む。従って原点の局所環も非被約である。
完備化の忠実平坦性とhull同定により、六項側のminiversal基底も非被約となる。
異なる次元の円錐成分は全て原点を通るので、この局所環は非純次元である。
斉次素イデアルによる商の原点完備化は整域である。実際、非零な二つの形式元の
最低次斉次項の積は、その斉次整域の中で非零となる。
従って相異なる極小素イデアルとそれらの次元は完備化後にも保たれ、hullも非純次元となる。
正則局所環の完全交叉商は純次元であるから、完全交叉ではない。
\end{proof}

三十八個の独立な二次生成元があることと、局所的な消去の後に五次の式が現れることは両立する。
前者は原点における埋め込みの最小生成次数、後者は非零Poisson点の近傍で
可逆な線形部分を消去した局所方程式の次数である。
また、同じ被約基底上の点集合が一致しても、上の冪零方向はArtin基底上の補完を変える。
本稿の全Artin基底上の節点条件と同様、補完問題ではイデアルの根基への置換は許されない。

\section{座標環による実現と補完定理の射程}

量子多項式環
\[
 S=k\langle x_1,x_2,x_3,w\rangle/
 (x_jx_i-q_{ij}x_ix_j,\ wx_i-\lambda_i x_iw),
 \qquad q_{ij},\lambda_i\ne0
\]
は、重み $(1,1,1,3)$ のHilbert級数
$((1-z)^3(1-z^3))^{-1}$ を持つ。
この場合はAuroux--Katzarkov--Orlovの定理をそのまま適用でき、
$S_{(0,1,3,4)}$ のperfect圏は $D^b(\operatorname{qgr}S)$ の許容部分圏となる。
別稿\cite{VKPoisson32}では、二変数四次式 $f$ に対し
\[
 \frac{k\langle y,z,x,w\rangle}{
 (yz-zy,\ xy-ayx,\ xz-bzx,\ wy-a^{-1}yw,\ wz-b^{-1}zw,\ wx-abxw+f)}
\]
という反復Ore拡大を構成する。
$a,b\ne0$ のとき、これはAS正則次元四、Gorenstein parameter六の4-Calabi--Yau代数である。
Orlovの定理\cite[Corollary 2.7]{Orlov2009}により標準六項のfull strong列を持つため、
この族についても実際の $\Perf(S_{(0,1,3,4)})$ の半直交実現が得られる。

一方、本稿の一般量子化定理が直接構成するのは、各次数成分が自由である形式的な重み付き代数と、
その六項の有限切り出しである。
任意の形式量子化に対するAS正則性、パラメータの代数化、あるいは解析的収束性はここでは使わない。
従って主定理の「周囲」は、明確に指定した $\Lambda$ の形式変形として理解する。
全ての非可換重み付き射影空間の分類という意味には広げない。

既存の比較で現れる
\[
 \operatorname{gr}_J B_{p,q}\simeq(Q_{p,q})_{(0,1,3,4)}
\]
は量子中間代数、すなわち随伴次数代数の実現を与える。
この等式だけから元の $B_{p,q}$ 自身の埋め込みは従わない。
本稿の $st=0$ 上のOre族では、随伴次数を取ることなく実際のcornerを同定している。

また、次数 $1,1,1,3$ の最小生成元と上のHilbert級数を持つ座標環では、
\[
 \dim S_3/(S_1S_2)=1,\qquad
 \operatorname{rank}(S_1\otimes S_2\longrightarrow S_3)=10
\]
である。したがって、この合成が階数十一である標準四項代数を、そのまま
$(0,1,3,4)$ の標準部分列と同定することはできない。
この階数制約は、標準生成対象を変更した導来埋め込みの不存在を意味しない。


\subsection{残る問題の区別}
Poissonスキームの少なくとも八成分と、三角形・共点三直線・共鳴点の局所環は別稿で決定した。
そのスキームは非被約である。しかし全成分の分類と最大次元は残る。
一般のPoisson量子化をnoetherian AS正則座標環で代数化する問題も、有限六項の形式補完とは別である。

別の例外列まで許した埋め込み問題には、本稿の標準補完障害を越える不変量が必要となる。
一方、同じ標準cornerに対する補完の同型類が高々一つであることは、
新しい幾何的実現を得たとき、その実現を既存の六項族と比較する道具となる。

\section{計算の再現範囲}

同梱のPython検証器は標準ライブラリの整数と有理数だけで実行できる。
通常実行は証明書を変更せず、再計算結果を保存済み証明書と比較する。
統合した入口は \texttt{code/verify\_v31.py} である。
Hochschild複体の微分、制限のchain map、評価行列の階数、Schouten式の全係数、
BVホモトピー、二変数族の全ての非恒等三重合成を検査する。

Poissonスキームの新規成分と局所環の検証は別稿の六つのゲートにまとめ、
\nolinkurl{code/verify_v32.py} から実行できる。

全Artin基底上の補完条件、全次数の量子化、二次的miniversal基底の同定は、有限計算の射程には含めない。
それぞれcorner同型の持上げ、Kontsevichの定理、二次生成とNakayamaの補題が証明の根拠である。

\begin{thebibliography}{99}
\interlinepenalty=10000
\bibitem{AKO}D.~Auroux, L.~Katzarkov and D.~Orlov,
\emph{Mirror symmetry for weighted projective planes and their noncommutative deformations},
Annals of Mathematics \textbf{167} (2008), 867--943.
\url{https://arxiv.org/abs/math/0404281}.

\bibitem{Beilinson}A.~A.~Beilinson,
\emph{Coherent sheaves on $\mathbb P^n$ and problems of linear algebra},
Functional Analysis and Its Applications \textbf{12} (1978), 214--216.
\url{https://doi.org/10.1007/BF01681436}.

\bibitem{DoninMakarLimanov1998}J.~Donin and L.~Makar-Limanov,
\emph{Quantization of quadratic Poisson brackets on a polynomial algebra of three variables},
Journal of Pure and Applied Algebra \textbf{129} (1998), 247--261.
\url{https://arxiv.org/abs/q-alg/9504002}.

\bibitem{FPSS}L.~Fu, M.~Porta, S.~Scherotzke and N.~Sibilla,
\emph{Hochschild--Kostant--Rosenberg isomorphism for derived Deligne--Mumford stacks},
arXiv:2509.00501; 著者公開版のCorollary~1.4とTheorem~4.15を参照。
\url{https://arxiv.org/abs/2509.00501}.

\bibitem{GG}J.~A.~Guccione and J.~J.~Guccione,
\emph{Hochschild cohomology of triangular matrix algebras},
arXiv:math/0104068v2 (2001).
\url{https://arxiv.org/abs/math/0104068}.

\bibitem{Kontsevich}M.~Kontsevich,
\emph{Deformation quantization of Poisson manifolds},
Letters in Mathematical Physics \textbf{66} (2003), 157--216.
\url{https://arxiv.org/abs/q-alg/9709040}.
最終版の\S6.3--6.4：
\url{https://www.ihes.fr/~maxim/TEXTS/DefQuant_final.pdf}.

\bibitem{Kuznetsov}A.~Kuznetsov,
\emph{Height of exceptional collections and Hochschild cohomology of quasiphantom categories},
Journal f\"ur die reine und angewandte Mathematik \textbf{708} (2015), 213--243.
\url{https://arxiv.org/abs/1211.4693}.

\bibitem{LiuXu1992}Z.-J.~Liu and P.~Xu,
\emph{On quadratic Poisson structures},
Letters in Mathematical Physics \textbf{26} (1992), 33--42.
\url{https://doi.org/10.1007/BF00420516}.

\bibitem{Pym}B.~Pym,
\emph{Quantum deformations of projective three-space},
Advances in Mathematics \textbf{281} (2015), 1216--1241.
\url{https://arxiv.org/abs/1403.6444}.

\bibitem{TWW}X.~Tang, X.~Wang and J.~J.~Zhang,
\emph{Twists of graded Poisson algebras and related properties},
Journal of Geometry and Physics \textbf{207} (2025), 105344.
\url{https://arxiv.org/abs/2206.05639}.

\bibitem{Orlov2009}D.~Orlov, \emph{Derived categories of coherent sheaves and triangulated categories of singularities},
in \emph{Algebra, Arithmetic, and Geometry: In Honor of Yu. I. Manin}, Vol.~II,
Progress in Mathematics \textbf{270}, Birkh\"auser, 2009, 503--531.
\url{https://doi.org/10.1007/978-0-8176-4747-6_16}.
\url{https://arxiv.org/abs/math/0503632}.

\bibitem{VKPoisson32}本研究、公表原稿IX、\emph{重み付き射影三次元空間のPoisson成分と非被約構造}, v32.
\nolinkurl{weighted_poisson_components_ja_v32.pdf}.
\end{thebibliography}

\end{document}
